\chapter{Concluzii și direcții de dezvoltare}

În lucrarea asta am proiectat și am construit CourseHub, o aplicație web pe
ASP.NET Core 8 prin care profesorii și elevii se găsesc și lucrează împreună.
Aplicația acoperă tot drumul unui utilizator: se înregistrează și își confirmă
contul prin OTP, își completează profilul, caută pe cineva potrivit, trimite o
cerere de colaborare și, după ce e acceptată, discută prin chat în timp real și
își trimit fișiere. În spate există un sistem real de roluri Identity, peste care
am construit panoul de administrare, iar toată interfața merge în română și în
engleză.

Sunt și lucruri pe care nu am apucat să le fac. Aplicația rulează pe o bază de
date locală, nu am publicat-o în cloud. Notificările de mesaje noi se iau prin
interogare periodică (\emph{polling}), nu instantaneu, iar un fișier trimis în
conversație se vede la celălalt abia după ce dă refresh. Nu am scris teste
automate. Iar dacă blochez un cont care e deja autentificat, sesiunea lui rămâne
activă până iese singur.

Dacă aș continua proiectul, aș lua-o în ordinea asta:

\begin{itemize}
    \item introducerea unei suite de teste automate, pentru a reduce riscul de
    regresie pe măsură ce platforma crește;
    \item invalidarea imediată a sesiunilor active la blocarea unui cont;
    \item extinderea livrării în timp real prin SignalR și pentru mesajele cu
    fișiere atașate, eliminând reîncărcarea paginii;
    \item publicarea platformei într-un mediu cloud, cu o bază de date gestionată
    și integrare continuă;
    \item organizarea materialelor partajate prin conversații într-o bibliotecă
    per colaborare;
    \item introducerea unui sistem de recomandare a partenerilor potriviți, bazat
    pe materiile de interes și pe istoricul colaborărilor.
\end{itemize}

În concluzie, CourseHub demonstrează fezabilitatea unei platforme dedicate
conectării profesorilor cu elevii, cu un accent particular pe securitate --
autorizare aplicată consecvent la nivel de server, validare riguroasă a
conținutului încărcat și moderare funcțională.
